% Part: lambda-calculus
% Chapter: introduction
% Section: reduction

\documentclass[../../../include/open-logic-section]{subfiles}

\begin{document}

\olfileid[pt]{lam}{int}{red}
\olsection{Redução de termos lambda}

O que se pode fazer com termos lambda? Simplificá-los. Se $M$ e $N$ são
termos lambda quaisquer e $x$ é uma variável qualquer, podemos usar
$\Subst{M}{N}{x}$ para denotar o resultado de tomar $N$ como substituto de~$x$
em~$M$, depois de renomear quaisquer variáveis ligadas de $M$ que, após a
substituição, interfeririam com as variáveis livres de~$N$. Por exemplo,
\[
\Subst{(\lambd[w][xxw])}{yyz}{x} = \lambd[w][(yyz)(yyz)w].
\]

\begin{digress}
Notações alternativas para substituição são $[N/x]M$, $[x/N]M$ e também
$M[x/N]$. Cuidado!
\end{digress}

Intuitivamente, $(\lambd[x][M])N$ e $\Subst{M}{N}{x}$ têm o mesmo
significado; o ato de substituir o primeiro termo pelo segundo é chamado de
\emph{$\beta$-contração}. $(\lambd[x][M])N$ é chamado de \emph{redex}, e
$\Subst{M}{N}{x}$, de seu \emph{contractum}. Em geral, se é possível
transformar um termo $P$ em~$P'$ por uma $\beta$-contração de algum
subtermo, dizemos que $P$ \emph{$\beta$-reduz-se a $P'$ em um passo} e
escrevemos $P \redone P'$. Se, a partir de $P$, podemos obter~$P'$ por algum
número de reduções em um passo (possivelmente nenhuma), então $P$
\emph{$\beta$-reduz-se} a~$P'$; em símbolos, $P \red P'$. Um termo que não
pode mais ser $\beta$-reduzido é chamado de \emph{$\beta$-irredutível}, ou
\emph{$\beta$-normal}. Diremos ``reduz-se'' em vez de
``$\beta$-reduz-se'' etc. quando o contexto for claro.

Consideremos alguns exemplos.
\begin{enumerate}
\item Temos
\begin{align*}
(\lambd[x][xxy]) \lambd[z][z] & \redone (\lambd[z][z])(\lambd[z][z]) y \\
& \redone (\lambd[z][z]) y \\
& \redone y.
\end{align*}
\item ``Simplificar'' um termo pode torná-lo mais complexo:
\begin{align*}
(\lambd[x][xxy])(\lambd[x][xxy]) & \redone (\lambd[x][xxy])(\lambd[x][xxy])y \\
& \redone (\lambd[x][xxy])(\lambd[x][xxy])yy \\
& \redone \dots
\end{align*}
\item A redução também pode deixar um termo inalterado:
\[
(\lambd[x][xx])(\lambd[x][xx]) \redone (\lambd[x][xx])(\lambd[x][xx]).
\]
\item Além disso, alguns termos podem ser reduzidos de mais de uma maneira;
  por exemplo,
\[
(\lambd[x][(\lambd[y][yx]) z)] v \redone (\lambd[y][yv]) z
\]
contraindo a aplicação mais externa; e
\[
(\lambd[x][(\lambd[y][yx]) z)] v \redone (\lambd[x][zx]) v
\]
contraindo a mais interna. Observe, porém, que nesse caso ambos os termos se
reduzem depois ao mesmo termo, $zv$.
\end{enumerate}

O resultado final no último exemplo não é coincidência; ele ilustra uma
propriedade profunda e importante do cálculo lambda, conhecida como
``propriedade de Church--Rosser''.

\end{document}
